% 中文版，用 xelatex 编译：xelatex main_zh.tex
\documentclass[10pt,UTF8,scheme=plain]{ctexart}
\usepackage[margin=1in]{geometry}
\usepackage{amsmath,amssymb}
\usepackage{booktabs}
\usepackage{enumitem}
\usepackage{xcolor}
\usepackage{caption}
\usepackage{tikz}
\usepackage[colorlinks=true,linkcolor=black,citecolor=blue!50!black,urlcolor=blue!50!black]{hyperref}
\newcommand{\T}[2]{#2}  % 图中文字：\T{English}{中文}
% Shared by the English and the Chinese build. Figure text is written as \T{English}{中文};
% each main file defines \T to pick one.
\usetikzlibrary{arrows.meta,positioning,calc,fit,backgrounds}
\definecolor{mainc}{HTML}{6B7280}   % mainline behaviour
\definecolor{offc}{HTML}{C2410C}    % off-mainline behaviour
\definecolor{wfc}{HTML}{15803D}     % workflow
\definecolor{verc}{HTML}{1D4ED8}    % verifier
\definecolor{hidc}{HTML}{9F1239}    % hidden from the system
\tikzset{
  >={Stealth[length=4pt,width=4pt]},
  card/.style={draw=#1, fill=#1!6, rounded corners=2pt, align=left, inner sep=4pt, font={\scriptsize\hyphenpenalty 10000 \exhyphenpenalty 10000 \relax}},
  card/.default=mainc,
  pill/.style={draw=#1, fill=#1!10, rounded corners=6pt, align=center, inner sep=3pt, font=\scriptsize},
  pill/.default=mainc,
  lab/.style={font=\scriptsize\bfseries, text=#1}, lab/.default=black,
  note/.style={font=\scriptsize\itshape, text=mainc, align=center},
}

\newcommand{\sys}{\textsc{ht2wv}}
\newcommand{\todo}[1]{\textcolor{red}{[#1]}}
\setlist{nosep,leftmargin=1.6em}
\captionsetup{font=small,labelfont=bf}
\renewcommand{\figurename}{图}
\renewcommand{\tablename}{表}
\renewcommand{\abstractname}{摘要}
\renewcommand{\refname}{参考文献}
\setlength{\parindent}{2em}
\linespread{1.15}

\title{\textbf{HumanTrace2Workflow\&Verifier}\\[2pt]
\large 从一个人已有的访问记录里编译出他的 workflow 与 verifier}
\author{Zhi Zheng\\ \small 新加坡国立大学 \\ \small\emph{工作草稿。结果是初步的：每种配置只跑了一次，样本为合成数据。}}
\date{}

\begin{document}
\maketitle

\begin{abstract}
人在电脑上重复着大量流程，但要把其中一个交给 agent，仍然得自己描述、自己写 prompt、自己检查结果。
最近的产品能把一次录下来的演示变成可复用的 skill，最近的研究能从被动记录的轨迹里归纳出任务模型。
两者恢复的都是一件事\emph{怎么做}，都没有恢复这个人认为怎样才算\emph{做对}。我们认为这两者不可分：
没有最优的 workflow，只有在某个人的验收标准之下最优的 workflow，所以要学的对象是一对 (workflow, verifier)。
verifier 的证据不在任务的主线上，而在主线之外：纠正、中断、复核、返工；其中很大一部分已经写在这个人
所用的 coding agent 的日志里。\sys{} 在逐项授权之下采集这类访问记录，确定性地把它们压缩成一条时间线
和一批“主线之外”的信号候选，然后在已有的 coding agent 里开无头会话，蒸馏出一份画像，并为每个任务蒸馏出
一个以多会话 prompt 形式交付的 workflow，以及一个分层的 verifier——每个会话拿到的都是同一份由它渲染成的 prompt。
为了度量，我们构造了合成样本：先预填一份个人风格，据此写出记录，再把风格藏起来；系统按两件事打分：
它的 workflow 满足了这个人的 rubric 中的几条，它的 verifier 抓住了几条被隐藏的风格。其中有一个持续七天、
带状态的样本，在里面把对的事按错的顺序做，结果就会不同。初步结果是：单步任务第一次执行就饱和了；
而学出来的 verifier 有一半会拒绝这个人本来会留下的产出——只有覆盖率、却不接受参考产出的，不能算 verifier。
三条边里只有“从人蒸馏 verifier”这一步会读到个人的原始记录；因此我们把这种样本格式看作训练本地开源模型的材料，
把许多不同的人的 (workflow, verifier) 看作一个多样化、能自我改进的 working group 的成员——那是本文之后的下一个阶段。
\end{abstract}

% ---------------------------------------------------------------------------------------------------------------
\section{引言}

Computer-use agent 已经能操作界面，但它\emph{该做什么}、结果要\emph{好到什么程度}，仍然由人手工指定。
有两条路线在减轻前一个负担：当前 agent 产品里的“录制并回放”功能把一次带讲解的演示变成 skill
\cite{recordskill,recordreplay,skillrecorder}；任务模型归纳则从被动记录的轨迹里恢复目标层级和过程\cite{tmi}。
两者产出的都是对流程的描述。后一个负担——判断结果能不能接受——仍然留给了人。

而这个标准是因人而异的。同样是写周报，两个人在“怎样才算合格”上可以完全不同：一个人在每个数字都对着源文件
核过、并在末尾写明数据来源之前，什么都不会发；另一个人要的是一段语气亲切的结论。为其中一个人优化出来的 agent，
对另一个人就是错的（图~\ref{fig:idea}a）。我们把这一点当真：

\begin{quote}
\emph{workflow 与 verification 是耦合的。要学的不是一个流程，而是一对 (workflow, verifier)，
而且 verifier 属于某一个人。}
\end{quote}

\paragraph{为什么要做这个系统：是三角，不是流水线。} 这里有三样东西，每一样作用于下一样（图~\ref{fig:triangle}）。
\emph{人指导 verifier}：他纠正了什么、拒绝了什么、回头核对了什么、在 rubric 上勾了什么，标准就是从这里来的；
没有他的批准，标准的任何改动都不生效。\emph{verifier 优化并加速 workflow}：它是搜索 workflow 时的目标函数；
同时它作为同一份 prompt 附给每一个 agent 会话，使一个很长的流程可以拆到多个会话里去做，而各个会话对
“做对了”的理解不会走散。\emph{workflow 给人提效}：它接手重复性的工作，并按这个人读得最快的方式汇报回来。
没有 verifier，人就只能直接去指导 workflow，而且每次运行都要指导一遍——那就是写 prompt，它不会累积。
有了 verifier，人的精力只在标准上花一次，此后这个循环每转一圈，都会留下更多记录，用来把标准磨得更准。

\begin{figure}[t]
\centering
\begin{tikzpicture}[x=1cm,y=1cm]
\node[card=mainc, align=center, text width=3.6cm, minimum height=1.2cm, font=\small] (H) at (8.25,5.7) {\textbf{\T{human}{人}}\\ {\scriptsize \T{purpose, standards, habits}{目的、标准、习惯}}};
\node[card=verc, align=center, text width=3.6cm, minimum height=1.2cm, font=\small] (V) at (3.7,1.45) {\textbf{verifier}\\ {\scriptsize \T{what counts as done, for this person}{对这个人来说，怎样才算做对}}};
\node[card=wfc, align=center, text width=3.6cm, minimum height=1.2cm, font=\small] (W) at (12.8,1.45) {\textbf{workflow}\\ {\scriptsize \T{a multi-session agent procedure}{由多个 agent 会话组成的流程}}};
\draw[->, line width=1.6pt, offc] (H) -- (V);
\draw[->, line width=1.6pt, verc] (V) -- (W);
\draw[->, line width=1.6pt, wfc] (W) -- (H);
\node[align=right, font=\scriptsize, text width=4.6cm, text=offc] at (2.95,4.05) {\textbf{\T{guides}{指导}}\\ \T{corrections, refusals and re-checks in their records; ticks the rubric; approves every change}{记录里的纠正、拒绝与复核；勾选 rubric；每一次改动都由本人批准}};
\node[align=center, font=\scriptsize, text width=7cm, text=verc] at (8.25,0.55) {\textbf{\T{optimises and accelerates}{优化并加速}}\\ \T{the objective the workflow is searched under, and one definition of done shared by all of its sessions}{既是搜索 workflow 时的目标函数，也是所有会话共用的同一份“做对了”的定义}};
\node[align=left, font=\scriptsize, text width=4.6cm, text=wfc] at (13.55,4.05) {\textbf{\T{boosts}{提效}}\\ \T{does the recurring work and reports back in the form this person reads fastest}{接手重复性的工作，并按这个人读得最快的方式汇报}};
\node[note, text width=4.4cm] at (8.25,2.75) {\T{each turn leaves more records,\\ and a sharper verifier}{每转一圈，记录更多，\\ verifier 更准}};
\node[draw=hidc, dashed, rounded corners=3pt, align=center, font=\tiny, text=hidc, text width=3.5cm, inner sep=3pt] at (2.2,5.75) {\T{the only step that reads raw personal records: keep it on a local, open-weight model}{唯一会读到个人原始记录的一步：放在本地的开源模型上做}};
\draw[hidc, dashed] (4.1,5.6) -- (6.8,4.42);
\end{tikzpicture}

\caption{为什么要有这个系统。人指导 verifier，verifier 优化并加速 workflow，workflow 把时间还给人。
三条边里只有第一条会接触到个人的原始记录。}
\label{fig:triangle}
\end{figure}

由此得到三个推论。个人化不是附加功能，而是问题定义本身。verifier 才是持久的资产：模型变强之后 workflow
可以重新生成，甚至换成完全不同的做法，而一个人接受什么变化得很慢。优化也因此有了明确的定义：
固定这个人的 verifier，在 workflow 的空间里搜索。

\paragraph{verifier 的证据在哪里。} 不在任务的主线上，而在 task mining 当作噪声丢掉的部分：撤销重做、
看一眼来源却什么都没复制、发送之前的迟疑；最密集的，是这个人怎么跟他的 coding agent 说话。每一句“这不对”
“你核对过吗”，每一次中断、每一次拒绝工具调用，都是他用自己的话说出来的验收标准，带时间戳，
被打回的产出就在上一条消息里（图~\ref{fig:idea}b）。这些日志本来就在机器上，不需要额外录制任何东西。

\begin{figure}[t]
\centering
\begin{tikzpicture}[x=1cm,y=1cm]
% ---------------- (a) one task, two people
\node[lab] at (3.15,5.3) {\T{(a) One task, two people}{(a) 同一个任务，两个人}};
\node[pill] (task) at (3.15,4.6) {\T{Task: write the weekly metrics report}{任务：写每周的指标周报}};
\node[card=verc, text width=2.75cm, minimum height=1.25cm] (va) at (1.55,3.25) {\textbf{\T{A accepts if}{A 的验收标准}}\\ \T{every number equals the source and the file is named}{每个数字都与源文件一致，并写明数据来源}};
\node[card=verc, text width=2.75cm, minimum height=1.25cm] (vb) at (4.75,3.25) {\textbf{\T{B accepts if}{B 的验收标准}}\\ \T{there is a short takeaway in a friendly tone}{有一段简短的结论，语气亲切}};
\node[card=wfc, text width=2.75cm, minimum height=1.25cm] (wa) at (1.55,1.45) {\textbf{\T{best for A}{对 A 最优的做法}}\\ \T{fetch, compute, re-check each number, then write}{取数、计算、逐个核对，再动笔}};
\node[card=wfc, text width=2.75cm, minimum height=1.25cm] (wb) at (4.75,1.45) {\textbf{\T{best for B}{对 B 最优的做法}}\\ \T{draft quickly, then polish the wording}{先快速起草，再打磨措辞}};
\draw[->] (task) -- (va.north); \draw[->] (task) -- (vb.north);
\draw[->, verc, thick] (va) -- (wa); \draw[->, verc, thick] (vb) -- (wb);
\node[note, text width=6cm] at (3.15,0.25) {\T{no best workflow, only the best one\\ under a given person's verifier}{没有最优的 workflow，\\ 只有某个人的 verifier 之下的最优}};
\draw[mainc!40] (6.45,0) -- (6.45,5.45);
% ---------------- (b) where the standard shows
\node[lab] at (11.5,5.3) {\T{(b) The standard shows beside the main line}{(b) 标准出现在主线之外}};
\foreach \px/\ptxt/\pcol in {7.35/{\T{asks for\\a draft}{让 agent\\起草}}/mainc, 9.0/{\T{agent\\drafts}{agent\\起草}}/mainc,
                      10.65/{\T{``numbers only,\\no praise''}{“只要数字，\\别夸”}}/offc, 12.3/{\T{``did you check\\the csv?''}{“你核对过\\csv 吗？”}}/offc,
                      13.95/{\T{refuses the\\send action}{拒绝\\发送动作}}/offc, 15.6/{\T{``ok''}{“可以”}}/mainc}
  \node[card=\pcol, align=center, text width=1.38cm, minimum height=0.9cm, inner sep=2pt] at (\px,4.1) {\ptxt};
\draw[->, mainc] (6.6,4.78) -- (16.45,4.78) node[above left, font=\tiny, text=mainc] {\T{time}{时间}};
\node[card=wfc, text width=3.7cm, align=center] (W) at (8.75,2.15) {\textbf{workflow}\T{: how it is done}{：怎么做}\\ \emph{\T{regenerated when models change}{模型换代后可以重新生成}}};
\node[card=verc, text width=3.7cm, align=center] (V) at (13.9,2.15) {\textbf{verifier}\T{: what counts as done}{：怎样算做对}\\ \emph{\T{kept, and refined with use}{长期保留，越用越准}}};
\foreach \px in {7.35,9.0,15.6} \draw[->, mainc!70] (\px,3.6) -- (W.north);
\foreach \px in {10.65,12.3,13.95} \draw[->, offc, thick] (\px,3.6) -- (V.north);
\draw[->, verc, thick] (V) -- node[above, font=\tiny, text=verc] {\T{objective}{目标函数}} (W);
\node[note, text width=9.6cm] at (11.5,0.55) {\T{grey: main-line events say how the work goes. orange: corrections, checks and\\ refusals, which task mining discards as noise, say what the person will not let pass}{灰色：主线上的事件说明事情怎么做。橙色：纠正、核对与拒绝，\\ 也就是 task mining 当作噪声丢掉的部分，说明这个人不放过什么}};
\end{tikzpicture}

\caption{立意。(a) 验收标准因人而异，而它决定了哪个 workflow 最优。(b) 这个标准的证据在任务主线之外，
在这个人与 coding agent 已有的对话里最密。}
\label{fig:idea}
\end{figure}

\paragraph{本文所处的位置。} 我们把“用 coding 模型把工作自动化”看成五个阶段（图~\ref{fig:ladder}）：
\emph{coding}，人在模型的帮助下写；\emph{coding agent}，一个 agent 把一件任务从头做到尾；
\emph{coding multi-agent workflow}，多个 agent 会话被编排成一个流程；
\emph{verifier-optimized coding multi-agent workflow}，这种编排不再由人手工完成，而是在一个明确的验收标准之下搜索出来；
以及\emph{automated (recursive self-improving) working group}，一个常设的小组，长期维护并改进它自己的 workflow。
每一个阶段都拿走一样原本要由人提供的东西：敲键盘、一步步盯着、编排步骤，最后是对编排本身的维护。
本文实现的是第四个阶段。让第四阶段成为可能的那个对象，也正是让第五阶段可以放心去尝试的那个对象：
一个明确的、可执行的、归本人所有的 verifier。有了它，被改进的东西（workflow）和评判它的东西（verifier）
就被分开了；一个会改写自己流程的小组，仍然要对一份它自己无权单方面改写的标准负责。

\paragraph{下一个阶段的关键是多样性。} 一个人的一对 (workflow, verifier)，只能造出一个照着这个人的样子工作的成员。
我们想要的 working group，是由从\emph{许多不同的人}那里蒸馏出来的许多对组成的：那个数字没核对就不发出去的分析师，
那个没有回归测试就不合并修复的工程师，那个按“谁先申请”来排队的售后运营。他们的 workflow 给小组带来一系列不同的做法；
他们的 verifier 给小组带来一系列不同的标准，这些标准之间的分歧是有用的——在一个人的标准下做出来的东西，
可以放到另一个人的标准下去审。这样的小组里的多样性不需要靠写 persona 的 prompt 凭空造出来：
它是从这些 (workflow, verifier) 背后的人那里继承来的，而且始终可以度量，因为每个成员的标准都是一个明确的、可执行的对象。

\paragraph{什么必须是私密的。} 图~\ref{fig:triangle} 的三条边里，只有一条会读到个人的原始记录：从人蒸馏出 verifier。
执行 workflow 看到的只是手头这件任务；优化 workflow 看到的是 verifier，而它是一段很短的、可以审阅的文本。
这两步都不需要这个人的历史。因此隐私要求只落在一个边界清楚的步骤上，而这一步就应该放在本人自己托管的模型上做。
下文介绍的那种样本——记录，配上生成这些记录的风格——正是这样的模型所需要的监督信号；所以我们正在收集的数据集，
同时也是训练能在本地运行的开源模型（Qwen 27B、35B-A3B 这个量级）的材料。更强的托管模型则可以用在执行和优化上，
那里不暴露任何私密的东西。

\begin{figure}[t]
\centering
\begin{tikzpicture}[x=1cm,y=1cm]
\foreach \lk/\lx/\ltitle/\lsup/\lcol in {
   1/1.55/{\T{coding}{写代码}}/{\T{writes the code,\\checks the result}{自己写代码，\\自己检查结果}}/mainc,
   2/4.85/{coding agent}/{\T{states the task,\\checks the result}{说明任务，\\检查结果}}/mainc,
   3/8.15/{\T{coding multi-agent\\workflow}{coding 多 agent\\workflow}}/{\T{arranges the sessions,\\checks the result}{编排各个会话，\\检查结果}}/mainc,
   4/11.45/{\T{verifier-optimised\\coding multi-agent\\workflow}{由 verifier 优化的\\coding 多 agent\\workflow}}/{\T{a standard, learned\\from their own records}{一份标准，\\从本人的记录里学出来}}/verc,
   5/14.85/{\T{automated, recursively\\self-improving\\working group}{自动化、递归自我\\改进的\\working group}}/{\T{pairs from many, diverse\\people; each approves\\changes to their standard}{来自许多不同的人的\\(workflow, verifier)；\\各自批准自己标准的改动}}/mainc}
{
  \node[card=\lcol, align=center, text width=2.75cm, minimum height=1.3cm] (L\lk) at (\lx,1.7) {\textbf{\ltitle}};
  \node[font=\tiny, text=mainc, align=center] at (\lx,0.45) {\lsup};
}
\foreach \lk [evaluate=\lk as \lnext using int(\lk+1)] in {1,...,4} \draw[->, thick] (L\lk.east) -- (L\lnext.west);
\node[draw=verc, very thick, rounded corners=3pt, fit=(L4), inner sep=2.5pt] {};
\node[font=\scriptsize\bfseries, text=verc] at (11.45,2.75) {\T{this paper}{本文}};
\node[font=\tiny\itshape, text=mainc] at (14.85,2.75) {\T{next stage}{下一阶段}};
\node[font=\tiny\itshape, text=mainc] at (8.2,-0.35) {\T{second row: what people still have to supply at each stage}{第二行：每个阶段仍然需要人来提供的东西}};
\end{tikzpicture}

\caption{我们设想的发展脉络。每个阶段都拿走一样原本要由人提供的东西。本文是第四阶段：多会话的 workflow
在一个从本人记录里学出来的 verifier 之下被优化。第五阶段把从许多不同的人那里蒸馏出来的 (workflow, verifier) 组成一个小组。}
\label{fig:ladder}
\end{figure}

\paragraph{贡献。}
\begin{itemize}
\item 把个人自动化表述为“人—verifier—workflow”的三角关系，其中要学的对象是一对 (workflow, verifier)；
      verifier 分五层（客观、顺序、风格、偏好、汇报），外加一份“这个人怎么查”的检查方式（\S\ref{sec:method}）。
\item \sys{}：一个开源实现。它在逐项授权之下读取一个人已有的访问记录；自己不调用任何模型 API，
      所有的阅读和写作都交给已有 coding agent 的无头会话去做，再用确定性的检查卡住这些会话的产出。
\item 一种带有隐藏的、预填的个人风格的样本格式与打分方式，其中包括一个对顺序敏感、持续七天、带状态的任务；
      以及在六个样本上的初步测量（\S\ref{sec:bench}--\ref{sec:results}）。
\end{itemize}

% ---------------------------------------------------------------------------------------------------------------
\section{问题}
\label{sec:problem}

一个人 $u$ 留下\emph{访问记录} $R_u$：与 coding agent 的对话、shell 历史、浏览器历史、导出的聊天记录、
文件名与修改时间，以及可选的屏幕录制。对一件反复出现的任务，这个人有一个\emph{目的}，并且说得出一份简短的
\emph{rubric} $Q_u$——他会逐条去勾的验收项。在它背后是一份更大的\emph{风格} $S_u$：他在用、却很少说出口的标准，
有的关于结果，有的关于做事的先后顺序，有的关于别人该怎么向他汇报。

系统从 $R_u$ 和 $Q_u$ 产出：
\begin{itemize}
\item 一个 \emph{workflow} $W$：一个有序的 prompt 列表，每个 prompt 对应一个 agent 会话；会话之间不共享记忆，
      只通过文件交接；
\item 一个 \emph{verifier} $V$：一组条目，每条是一个判断，配有可执行的检查脚本或评审用的 rubric，带有所属的层、权重，
      以及指向支撑它的记录的指针；外加一份\emph{检查方式}（policy），描述这个人怎么查：查多深、什么时候查、
      失败了怎么办、什么情况下停下来问人。
\end{itemize}
$V$ 被渲染成一份 prompt，$W$ 的每一个会话都会拿到它。于是流程可以切成任意多个会话，但“做对了”只有一份定义。

在一个留出的任务实例上（其参考产出为 $y^\star$，即 $u$ 自己会留下的那一份），我们度量三件事：
\begin{align*}
\textbf{rubric 完成度} &:\quad |\{q \in Q_u : q(W(x)) \text{ 成立}\}| \,/\, |Q_u|, \\
\textbf{风格覆盖率} &:\quad |\{s \in S_u : V \text{ 在某个对 } y^\star \text{ 判通过的条目上判 } y^\star_{\neg s} \text{ 不通过}\}| \,/\, |S_u|,\\
\textbf{参考产出被接受} &:\quad V(y^\star) = \text{accept},
\end{align*}
其中 $y^\star_{\neg s}$ 是把参考产出恰好按 $s$ 所禁止的方式改坏之后的版本。覆盖率奖励严格的 verifier；
“参考产出被接受”则惩罚那些在错误的地方严格的 verifier。对长程任务，rubric 完成度按天报告，
主要的数字是有几天被\emph{填满}，即当天适用的 rubric 条目全部通过。

% ---------------------------------------------------------------------------------------------------------------
\section{一个例子：一周的退款}
\label{sec:example}

图~\ref{fig:week} 是我们最难的一个样本。许晴在一家小店做售后退款。每天她收到一张退款申请表，
在一个后台里逐条处理，而后台里的每个动作都是真的：退出去的钱收不回来，发给买家的通知也收不回来。
后台每天最多允许退 \textyen 1000，超出的会被拒绝。

给系统的东西很少：目的；一份七条的 rubric（结果要和她自己做的一样；每一单在动手之前都先查询；
通知买家是最后一步且只通知一次；没有被拒绝的退款、没有退两次的订单；汇报的第一行给出当天的总数；周日给一份周汇总）；
以及她的访问记录——大约六十轮 agent 对话、二十三条 shell 命令、二十三条聊天消息和五条浏览记录，其中约一半与任务无关。

\emph{没有}给系统的，是写出这些记录时所依据的十七条标准。六条管当天之内的顺序与前置条件
（先申请的先处理，而不是按表格的行顺序；买家自己原因的退货要扣运费；超过 \textyen 500 的升级；
先退款再通知；退货没签收不退；不为了凑额度把一笔退款拆开）。五条管她希望别人怎么向她汇报
（第一行就是结论；一单一行，不写过程；需要她处理的放最前面并标出来；金额写成 \textyen{} 加两位小数；不贴日志）。
还有六条只在跨天时才存在（昨天没退成的单，今天按它最初的申请时间和新申请排成一队；主管批准的单，批准当天就退；
寄回并签收之后的再次申请不算重复；已经退过款之后的再次申请才算；被驳回的要通知买家；周汇总遵守同样的汇报规矩）。

\begin{figure}[t]
\centering
\begin{tikzpicture}[x=1cm,y=1cm]
% visible / hidden
\node[card=mainc, text width=7.4cm] (vis) at (4.0,7.75) {\textbf{\T{Given to the system}{给系统看的}}\\
\T{purpose: ``handle each day's refunds as I would, then report to me''}{目的：“每天的退款按我的做法处理完，再向我汇报”}\\
\T{rubric R1--R7: the items she will tick}{rubric R1--R7：她会逐条勾的验收项}\\
\T{records: agent logs, shell history, chats, browsing}{记录：agent 对话、shell 历史、聊天、浏览}};
\node[card=hidc, text width=7.4cm] (hid) at (12.35,7.75) {\textbf{\T{Withheld: the pre-filled style the records were written from}{隐藏的：预先填好、用来生成这些记录的个人风格}}\\
\T{6 in-day ordering rules}{6 条当天内的顺序规则} \hfill S1--S5, S10\\
\T{5 reporting-style rules}{5 条汇报风格} \hfill S6--S9, S11\\
\T{6 cross-day rules}{6 条跨天规则} \hfill L1--L6};
% days
\foreach \dk/\dn/\dtxt in {0/{\T{Mon}{周一}}/{\T{8 requests\\cap reached\\W107 waits}{8 条申请\\额度用完\\W107 等明天}},
                      1/{\T{Tue}{周二}}/{\T{W102 approved\\W107 goes first\\2 orders wait}{W102 获批\\W107 排最前\\2 单等明天}},
                      2/{\T{Wed}{周三}}/{\T{W104 returned\\W101 repeats\\W111 over 500}{W104 已寄回\\W101 重复申请\\W111 超 500}},
                      3/{\T{Thu}{周四}}/{\T{W111 rejected\\cap reached\\W116 waits}{W111 被驳回\\额度用完\\W116 等明天}},
                      4/{\T{Fri}{周五}}/{\T{W116 goes first\\W118 over 500\\W121 waits}{W116 排最前\\W118 超 500\\W121 等明天}},
                      5/{\T{Sat}{周六}}/{\T{W118 approved\\W113 returned\\W121 waits}{W118 获批\\W113 已寄回\\W121 继续等}},
                      6/{\T{Sun}{周日}}/{\T{W121 goes first\\W125 over 500\\\textbf{week report}}{W121 排最前\\W125 超 500\\\textbf{周汇总}}}}
  \node[card=mainc, text width=1.92cm, minimum height=1.55cm, inner sep=3pt] (day\dk) at (\dk*2.345+1.14,5.45) {\textbf{\dn}\\[1pt] \dtxt};
\foreach \dk [evaluate=\dk as \dnext using int(\dk+1)] in {0,...,5}
  \draw[->, offc, very thick] (day\dk.east) -- (day\dnext.west);
\node[note, text=offc, text width=15.5cm] at (8.18,4.3) {\T{arrows: the back-office state carries over. the cap resets every day; whatever was left undone is still there tomorrow}{箭头：后台状态延续到第二天。额度每天重置；没办完的事明天还在}};
% order matters
\node[lab] at (8.18,3.65) {\T{Why order matters: the same Monday, cap \textyen 1000}{为什么顺序要紧：同一个周一，额度 \textyen 1000}};
\node[card=wfc, text width=7.4cm] (good) at (4.0,2.2) {\textbf{\T{her order: oldest request first}{她的顺序：先申请的先处理}}\\
08:50 W103 \textyen 259 \quad 09:02 W101 \textyen 187\\
09:10 W105 \textyen 150 \quad 09:15 W106 \textyen 329 \quad $\Rightarrow$ \T{\textyen 925 refunded}{共退 \textyen 925}\\
\T{09:50 W107 \textyen 210 no longer fits: \emph{the latest request waits}}{09:50 的 W107 \textyen 210 放不下：\emph{最晚申请的等明天}}};
\node[card=offc, text width=7.4cm] (bad) at (12.35,2.2) {\textbf{\T{sheet order: rows as they come}{表格顺序：按行处理}}\\
W101 \textyen 187 \quad W107 \textyen 210\\
W106 \textyen 329 \quad W103 \textyen 259 \quad $\Rightarrow$ \T{\textyen 985 refunded}{共退 \textyen 985}\\
\T{W105 \textyen 150 no longer fits: \emph{09:10 loses to 09:50}}{W105 \textyen 150 放不下：\emph{09:10 申请的输给了 09:50 的}}};
\node[note, text width=15.5cm] at (8.18,0.85) {\T{every single action is valid in both columns. only the order differs, and the back office ends in a different state}{两边的每一个动作单看都合法，只有顺序不同，后台的最终状态却不一样}};
\end{tikzpicture}

\caption{七天的样本。上：系统看得到的和被隐藏的。中：这一周；没退成的单、主管的审批、寄回的包裹把各天连在一起。
下：同一个周一，按两种顺序处理。}
\label{fig:week}
\end{figure}

这个任务是特意设计成“看起来正确、但顺序不对的 workflow 会失败”的。图~\ref{fig:week} 的下半部分是周一。
按申请时间处理，四单被退款，最晚申请的那一单等明天的额度。按表格的行顺序处理，每一个动作单看仍然是合法的，
但一个较晚的申请先占掉了额度，一个较早的申请反而被挡在外面；后台最终停在另一个状态，到了周二，排在队首的也成了另一单。
参考的一周和每一个“顺序类”探针，都是真的把后台跑一遍得到的：一次完全照她的规则，然后每条规则各一次、
只忽略这一条。这样日志、状态和汇报三者始终互相一致——verifier 只有知道那条规则才抓得住探针，
而不能靠发现文件之间对不上。

% ---------------------------------------------------------------------------------------------------------------
\section{方法}
\label{sec:method}

\begin{figure}[t]
\centering
\resizebox{\textwidth}{!}{\begin{tikzpicture}[x=1cm,y=1cm]
% row 1: pipeline
\node[card=mainc, text width=2.45cm, align=center] (rec) at (1.4,6.3) {\textbf{\T{access records}{访问记录}}\\ \T{agent logs, shell and browser history, chats}{agent 对话、shell 与浏览器历史、聊天}\\ \emph{\T{per-source consent}{逐项授权}}};
\node[card=mainc, text width=2.45cm, align=center] (cmp) at (4.85,6.3) {\textbf{\T{compile}{编译}}\T{, no model}{（不调模型）}\\ \T{timeline plus off-mainline signal candidates}{时间线 + 主线之外的信号候选}};
\node[card=mainc, text width=2.45cm, align=center] (dst) at (8.3,6.3) {\textbf{\T{distil}{蒸馏}}\T{, agent sessions}{（agent 会话）}\\ \T{A: profile and tasks}{A：画像与任务}\\ \T{B: one skill per task}{B：每个任务一个 skill}};
\node[card=wfc, text width=4.2cm, align=center] (wf) at (13.1,7.05) {\textbf{workflow}\\ \T{stage prompts 01, 02, 03}{各阶段的 prompt：01、02、03}};
\node[card=verc, text width=4.2cm, align=center] (vf) at (13.1,5.55) {\textbf{verifier}\\ \T{objective, process, style,\\ preference, report $+$ policy}{客观、顺序、风格、偏好、汇报\\ $+$ 检查方式}};
\draw[->] (rec) -- (cmp); \draw[->] (cmp) -- (dst);
\draw[->] (dst.east) -- (wf.west); \draw[->] (dst.east) -- (vf.west);
\node[font=\tiny, text=mainc, align=center] at (10.35,6.3) {\T{must\\pass lint}{须通过\\lint}};
% row 2: execution
\node[lab] at (1.4,3.2) {\T{execute}{执行} $\rightarrow$};
\foreach \sk/\sx in {1/3.9,2/6.9,3/9.9}
  \node[card=wfc, text width=2.2cm, align=center] (s\sk) at (\sx,3.2) {\textbf{\T{session}{会话} \sk}\\ \T{stage 0\sk{} prompt}{阶段 0\sk{} 的 prompt}\\ \textcolor{verc}{$+$ \T{verifier prompt}{verifier 的 prompt}}};
\node[card=mainc, text width=2.2cm, align=center] (out) at (13.1,3.2) {\textbf{\T{output}{产出}}\\ \T{and a report\\ to the person}{以及给本人的\\ 汇报}};
\draw[->] (s1) -- node[above, font=\tiny] {\T{files}{文件}} (s2); \draw[->] (s2) -- node[above, font=\tiny] {\T{files}{文件}} (s3); \draw[->] (s3) -- (out);
\draw[verc, thick] (vf.south) -- (13.1,4.5) -- (3.9,4.5);
\foreach \sx in {3.9,6.9,9.9} \draw[->, verc, thick] (\sx,4.5) -- (\sx,3.78);
\node[font=\tiny, text=verc] at (8.0,4.68) {\T{one shared definition of done, attached to every session}{同一份“做对了”的定义，附给每一个会话}};
% row 3: signals
\node[card=mainc, text width=2.6cm, align=center] (fb) at (13.1,1.1) {\T{she ticks the \textbf{rubric}\\ (stated, visible)}{本人勾 \textbf{rubric}\\ （说出口的、可见的）}};
\node[card=mainc, text width=3.2cm, align=center] (it) at (8.0,1.1) {\T{\textbf{iterate} session\\ repair stages; add verifier\\ items backed by evidence}{\textbf{迭代}会话\\ 修各阶段；补上有证据的\\ 验收条目}};
\node[card=hidc, text width=1.55cm, align=center] (pr) at (15.75,3.2) {\textbf{\T{hidden style}{隐藏风格}}\\ \T{probes give coverage}{探针给出覆盖率}};
\draw[->] (out) -- (fb);
\draw[->] (fb) -- (it);
\draw[->] (it.west) -- (-0.3,1.1) -- (-0.3,8.0) -- (13.1,8.0) -- (wf.north);
\node[font=\tiny, text=mainc, rotate=90] at (-0.5,4.5) {\T{revised skill}{修订后的 skill}};
\draw[->, hidc, dashed, thick] (pr.north) -- (15.75,5.55) -- (vf.east);
\node[font=\tiny, text=hidc, align=center] at (15.75,2.2) {\T{evaluation only,\\never fed back}{只用于评测，\\从不反馈}};
\end{tikzpicture}
}
\caption{流程。记录在授权之下采集，并在不调模型的情况下被压缩；agent 会话蒸馏出画像，并为每个任务蒸馏出
workflow 和 verifier；workflow 以多个会话的方式执行，每个会话拿到的都是同一份 verifier prompt。
rubric 上的失败驱动修订；隐藏的风格条目只用于打分。}
\label{fig:method}
\end{figure}

\paragraph{采集：先授权。} 每个来源在被发现时都不会被打开，只有本人允许之后才读取：本机和服务器上的 coding agent 日志、
shell 历史、浏览器历史（只保留域名和路径）、导出的聊天记录、最近打开的文件、git 提交。屏幕录制和录屏视频导入是可选项。
所有文本入库前都会脱敏。每个来源写的都是同一种事件格式。

\paragraph{编译：不调模型。} 一天的事件被压成一条时间线和一批\emph{主线之外的信号候选}。对话类信号只需要日志：
agent 回复之后的一次纠正、一次要求给出证明、一次被中断的轮次、一次被拒绝的工具调用、一份用户自己写的任务规格。
屏幕类信号需要录屏：看一眼另一个窗口却什么都没复制回来（这是核对；复制回来的则是取数）、连续撤销、
人已经走开之后又回来重新保存的文件（这不需要任何标注，就给出了同一份产物的“被丢弃版”和“被保留版”）。
检测器故意做成高召回的；每个命中只是一个留给 agent 去解读的指针，不是结论。

\paragraph{蒸馏：在别人的 agent 里做。} \sys{} 不调用模型 API。它写出一份任务书（产出、硬性约束、必须停下的情况、自检），
然后在已安装的 coding agent 里开一个无头会话。一个会话写画像——这个人的工作习惯、验收风格、汇报风格、做事顺序，
每条都带证据——并列出候选任务。然后每个任务一个会话，写出 skill。

\paragraph{这一对。} workflow 是一串阶段 prompt。verifier 有五层：\emph{客观}（与人无关的正确性）、
\emph{顺序}（动作的先后与前置条件，在执行留下的日志上检查）、\emph{风格}（这个人，跨任务）、
\emph{偏好}（这个人，这个任务）和\emph{汇报}（汇报回来的东西的结构、长度与可读性）。
每一条个人条目都必须引用记录，否则只能标成低置信度的假设。lint 这一步会拒绝这样的 verifier：
它判这个人留下的产出不通过，或者它没有任何条目能判这个人丢弃的产出不通过。

\paragraph{执行与修订。} 每个阶段作为一个独立的会话运行；verifier 的 prompt 附在所有会话后面，最后一个阶段负责跑可执行的检查。
一次运行之后，这个人在 rubric 上的结论交给一个\emph{迭代}会话，它可以修各个阶段，也可以补上有证据支撑的验收条目。
在正式的优化循环里，修订 workflow 时 verifier 是按哈希冻结的，负责修订的会话看不到留出集上的结果；
而 verifier 的任何改动，都要等本人批准之后才生效。

% ---------------------------------------------------------------------------------------------------------------
\section{带隐藏风格的样本}
\label{sec:bench}

一个样本是按“一个人存在的顺序”来构造的：先有风格，再有行为。我们先写人物、目的、rubric 和一份风格清单；
然后写访问记录，让每一条风格留下它自然会留下的那种痕迹（一份被打回的草稿、一句“你核对过吗”、对同事说的一句话、
一次被拒绝的动作），并混入与任务无关的日常活动。记录被渲染成各个来源的原生格式（agent 会话日志、带时间戳的 shell 历史、
浏览器的历史数据库、聊天导出），再由真实的采集器读回来。风格清单和参考产出被隐藏。
有效性检查要求：每一条风格至少影响了一条记录，并且参考产出能通过样本自己的 rubric。表~\ref{tab:samples} 列出了六个样本。

\begin{table}[h]
\centering\footnotesize\setlength{\tabcolsep}{4pt}
\begin{tabular}{@{}llllcc@{}}
\toprule
样本 & 人物 & 任务 & 记录来源 & 隐藏条目 & 时间跨度 \\
\midrule
weekly-report & 增长分析师 & 两份 csv $\rightarrow$ 周报 & 2 个 agent、shell、浏览、聊天 & 7 & 单步 \\
swe-bugfix & 后端工程师 & issue $\rightarrow$ 修复、测试、提交信息 & agent、shell、浏览、聊天 & 7 & 单步 \\
guest-reply & 民宿老板 & 客人消息 $\rightarrow$ 回复 & agent、浏览、聊天 & 7 & 单步 \\
video-publish & 美食视频博主 & 拍摄笔记 $\rightarrow$ 发布信息 & agent、浏览、聊天 & 7 & 单步 \\
refund-ops & 售后运营 & 一天的退款，带状态 & agent、shell、浏览、聊天 & 11 & 1 天 \\
refund-week & 售后运营 & 七天的退款，带状态 & agent、shell、浏览、聊天 & 17 & 7 天 \\
\bottomrule
\end{tabular}
\caption{样本。在两个退款样本里，动作的顺序会改变最终状态。}
\label{tab:samples}
\end{table}

% ---------------------------------------------------------------------------------------------------------------
\section{初步结果}
\label{sec:results}

\paragraph{设置。} 所有会话都是 Codex CLI 会话，模型为 \texttt{gpt-6-astra}，推理强度 \texttt{xhigh}，在一台服务器上运行。
每个样本只跑了一次。对单步样本，我们先执行建好的 skill，把 rubric 上的结论交给一个迭代会话，再执行一次；
需要评审判断的 verifier 条目只在最后一轮评估。

\begin{table}[h]
\centering\footnotesize\setlength{\tabcolsep}{4pt}
\begin{tabular}{@{}lccccc@{}}
\toprule
 & \multicolumn{2}{c}{rubric 完成度} & \multicolumn{2}{c}{风格覆盖率} & 参考产出 \\
\cmidrule(lr){2-3}\cmidrule(lr){4-5}
样本 & 第 0 轮 & 第 1 轮 & 只算检查脚本 & 加上评审条目 & 被接受 \\
\midrule
weekly-report & 5/5 & 5/5$^{\dagger}$ & 6/6 & 6/6 & 是 \\
swe-bugfix    & 4/5 & 5/5 & 1/6 & 6/6 & \textbf{否} \\
guest-reply   & 5/5 & 4/5 & 4/6 & 6/6 & \textbf{否} \\
video-publish & 5/5 & 5/5 & 5/6 & 6/6 & 是 \\
\bottomrule
\end{tabular}
\caption{单步样本，各跑一次。每个样本有一条隐藏条目是过程类标准，没有针对产出的探针，不计分。
$^{\dagger}$最初被判为 2/5，原因是我们的一条 rubric 错误地要求列表符号；表中是修正 rubric 之后对同一份产出重新打的分。}
\label{tab:single}
\end{table}

\paragraph{单步任务太简单了。} 四个 workflow 里有三个在第一次执行时就满足了整份 rubric；把评审条目算上之后，
所有带探针的隐藏条目都被抓住了（表~\ref{tab:single}）。在这些样本里，产出是一个静态文件，做的顺序无关紧要，
每条标准又几乎被本人在某次对话里原话说过。它们验证了管线是通的，除此之外说明不了多少——这也是我们去做退款样本的原因。

\paragraph{只看覆盖率会被误导。} 四个样本里有两个，经过一次修订之后的 verifier \emph{拒绝了参考产出}——
也就是这个人本来会留下的那份回复或补丁。在 \emph{guest-reply} 里，修订后的 verifier 在三条硬性条目上判参考回复不通过：
提前入住的加收费用没有写明“视当天房态”这个条件；报出的总价没有包含宠物清洁费；产出里没有能证明“这条回复没有被发送”的执行记录。
每一条单独看都说得通，但没有一条是她的标准。同一次修订还让 workflow 变得过于谨慎：它在一份正确的回复旁边
又写了一份“已阻塞”的说明，丢掉了一条之前已经满足的 rubric。一个什么都判不通过的 verifier 能拿到满分的覆盖率；
所以“参考产出被接受”必须和覆盖率一起报告，而修订这一步需要拿这个人留下的产出当回归集。

\paragraph{打分本身也很容易引入假象。} \emph{weekly-report} 看上去从 5/5 掉到 2/5，其实是我们自己的问题：
rubric 要求了这个人从没要求过的列表符号。我们把这件事留在表里，因为它和上一段是同一种失败，只不过发生在评测这一侧：
一个过于具体的检查，冒充了一条标准。

\paragraph{那一周。} \todo{refund-week 的结果：七天里填满了几天、逐天的 rubric 矩阵、按组的风格覆盖率
（当天顺序 / 汇报 / 跨天）、逐天的参考产出是否被接受。写作时运行尚未结束；从 \texttt{bench-report.md} 填入。}

\paragraph{需要说明的限制。} 这些数字来自一次运行、一个模型，以及由我们自己为了体现风格而写出来的记录，
所以它们更接近上界，而不是估计值。目前还没有真人参与。系统还没有在真人的记录上、带着 agent 后端完整地跑过一遍；
在作者自己的机器上，只有确定性的那几个阶段在真实日志上运行过。

% ---------------------------------------------------------------------------------------------------------------
\section{相关工作}
\label{sec:related}

\emph{从演示到 skill。} 现在的产品能把一次录下来的、有时带讲解的演示变成可复用的 skill
\cite{recordskill,recordreplay,skillrecorder}。它们是一次性的，需要人专门演示一遍，记录的是主线。
\emph{从被动轨迹到任务模型与用户模型。} 任务模型归纳\cite{tmi}能把未标注轨迹里交错的任务拆开，并构建带控制流的目标层级；
通用用户模型\cite{gum}从截图里维护关于用户的、带置信度的命题。两者都不产出验收标准。
一项从 GUI 轨迹里挖掘 skill 文件的研究报告说，聚出来的簇可读，但对下游几乎没有帮助\cite{skillmd}；
我们把它读作“缺一个目标函数”。\emph{从 agent 轨迹里得到 workflow}\cite{awm}，其“成功”由基准的评测器给出。
\emph{带 verifier 的 skill。} EvoSkills\cite{evoskills} 让 skill 生成器与一个替代 verifier 共同演化，
并报告去掉 verifier 后损失很大；agentic rubrics\cite{rubrics} 让 agent 读完仓库后写出一份带权重的检查清单。
它们的 verifier 是从任务推出来的，不是从某个人那里来的。\emph{从隐式反馈里学偏好。}
PRELUDE\cite{prelude} 从用户对 agent 产出的编辑里推断潜在偏好，并用来约束生成。
我们把信号从“编辑”扩展到纠正、拒绝和做事顺序上的习惯，并把结果用作一个可执行的验收测试。
\todo{本节依据的是摘要和二手综述；投稿前需要逐条对照原文核实。}

% ---------------------------------------------------------------------------------------------------------------
\section{局限与下一步}

多样化的 working group 和本地蒸馏目前都还不存在：没有任何开源模型在这些样本上被训练或评测过，
本文报告的所有会话都运行在托管模型上。核心主张——在别人的 verifier 之下优化出来的 workflow，在本人这里被接受得更少——
还没有被检验；这需要两个人（或者同一任务上的两种风格）和一次交叉评测。合成的记录比真实的记录友好：
大多数标准都被本人直接说了出来，也很少有东西与之矛盾。更难的样本应当包含：只被做出来、从没被说出来的标准；
随时间变化的标准；以及习惯不同、却出现在同一个聊天里的同事（退款样本里已经有几处）。
对留出材料的隔离靠的是工作目录，而不是操作系统级的沙盒。最后，verifier 最多只能和这个人自己的盲点一样好；
客观层的存在，以及每一条条目都对本人可见，是目前对这一点的回答。

\begin{thebibliography}{9}\small
\bibitem{recordskill} Anthropic. Record a Skill in Claude Cowork. Product feature, 2026.
\bibitem{recordreplay} OpenAI. Record \& Replay in the Codex desktop app. Product feature, 2026.
\bibitem{skillrecorder} Microsoft. Skill Recorder. \url{https://github.com/microsoft/skill-recorder}, 2026.
\bibitem{tmi} Y.~Jiang, Z.~Z.~Wang, R.~Chen, D.~Yang. Inducing Task Models from Computer-Use Traces. arXiv:2608.20319, 2026.
\bibitem{gum} O.~Shaikh et al. Creating General User Models from Computer Use. arXiv:2505.10831, 2025.
\bibitem{skillmd} Mining skill files from GUI trajectories. arXiv:2606.20363, 2026. \todo{作者、准确标题}
\bibitem{awm} Z.~Z.~Wang et al. Agent Workflow Memory. arXiv:2409.07429, 2024.
\bibitem{evoskills} EvoSkills: co-evolving skills and surrogate verifiers. arXiv:2604.01687, 2026. \todo{作者、准确标题}
\bibitem{rubrics} Agentic Rubrics as Contextual Verifiers for SWE Agents. arXiv:2601.04171, 2026. \todo{作者}
\bibitem{prelude} G.~Gao, A.~Taymanov, E.~Salinas, P.~Mineiro, D.~Misra. Aligning LLM Agents by Learning Latent Preference from User Edits. NeurIPS, 2024.
\end{thebibliography}

\end{document}
